\documentclass[11pt,reqno]{amsart}
\usepackage[T1]{fontenc}
\usepackage{amsmath,amssymb,amsthm,mathtools}
\usepackage[margin=1in]{geometry}
\usepackage[colorlinks=true,linkcolor=blue,citecolor=blue,urlcolor=blue]{hyperref}
\setlength{\emergencystretch}{2em}
\newtheorem{theorem}{Theorem}[section]
\newtheorem{proposition}[theorem]{Proposition}
\newtheorem{lemma}[theorem]{Lemma}
\theoremstyle{remark}\newtheorem{remark}[theorem]{Remark}
\newcommand{\R}{\mathbb R}
\title[Kravchuk Bilinear Formula and the Mehler Limit]{From the Kravchuk Bilinear Formula\\to the Mehler Kernel}
\date{A focused development}
\begin{document}
\begin{abstract}
We establish the finite Kravchuk bilinear formula from a binary
independent-site model and a hypergeometric summation identity. We
then use the closed finite formula, in an equivalent coefficient form,
to prove its central scaling limit by a one-dimensional saddle-point
calculation. The result is the Mehler kernel. Probability is used to
identify the finite kernel and its normalization, not as a substitute
for the asymptotic calculation. No interchange of a growing polynomial
sum with a limit is assumed.
\end{abstract}
\maketitle
\section*{Plan and conventions}
The argument has one direction:
\[
\text{binary occupation model}
\ \longrightarrow\ \text{finite bilinear formula}
\ \longrightarrow\ \text{coefficient asymptotics}
\ \longrightarrow\ \text{Mehler kernel}.
\]
The general induced-matrix viewpoint is described in
\cite{FeinsilverSFAII2023}; cumulative trials provide a concrete
sampling mechanism \cite{HoareRahman1984}. We retain only the binary
construction needed here. Throughout $N\ge1$ is an integer,
$0\le k,\ell\le N$, and the correlation parameter satisfies $0<r<1$.
The central coordinate and its mesh size are
\begin{equation}\label{eq:coordinates}
x_{k,N}=\frac{N-2k}{\sqrt{2N}},\qquad
\delta_N=\sqrt{\frac2N}.
\end{equation}
The limiting probability measure is
$d\mu(x)=g(x)dx$, where $g(x)=\pi^{-1/2}e^{-x^2}$; its variance is $1/2$.

\section{The finite probabilistic basis}
\subsection*{Cumulative Bernoulli trials: save, then retry}
We first specify what one update of the process means. There are $N$
labelled sites, each either a success (state $1$) or a failure (state
$0$). Suppose the input contains $k$ successes. Choose parameters
$0\le\alpha<1$ and $0<\beta<1$.
\begin{enumerate}
\item \emph{Save successes.} Independently save each of the $k$ current
successes with probability $\alpha$. A failure cannot be saved.
If $J$ successes are saved, then $J\sim\operatorname{Bin}(k,\alpha)$.
\item \emph{Retry all other sites.} Leave the $J$ saved successes
unchanged. Each of the other $N-J$ sites receives a fresh Bernoulli
trial with success probability $\beta$, independently of the saving
stage and of all other fresh trials. This includes both old failures
and old successes that were not saved.
\item \emph{Count the output.} If $R$ fresh trials succeed, the new
success count is $L=J+R$, where
$R\mid J=j\sim\operatorname{Bin}(N-j,\beta)$.
\end{enumerate}
Thus saved successes and newly obtained successes accumulate within
one update. The count need not increase: an old success that is not
saved can fail its fresh trial. On the next update the same rule is
applied anew; no success is permanently protected.

Writing $B_{m,p}(h)=\binom mhp^h(1-p)^{m-h}$, with value zero outside
$0\le h\le m$, the rule gives the explicit positive transition sum
\[
\Pr(L=\ell\mid k\text{ input successes})
=\sum_{j=0}^{\min(k,\ell)}B_{k,\alpha}(j)B_{N-j,\beta}(\ell-j).
\]
This is the saved-and-recouped-successes interpretation used here,
following the cumulative-trial viewpoint of \cite{HoareRahman1984}.
The independent one-site description below is the same update, not a
second probabilistic model.

\subsection{One-site transition and occupations}
For an input failure, the output is a success with probability
$\beta$. For an input success, it is a success either by being saved,
or by not being saved and then succeeding on the fresh trial.
The one-site matrix, with rows denoting input states and columns output
states in the order $0,1$, is therefore
\[
P_{\alpha,\beta}=
\begin{pmatrix}
1-\beta&\beta\\
(1-\alpha)(1-\beta)&\alpha+(1-\alpha)\beta
\end{pmatrix}.
\]
To obtain the symmetric kernel needed for the limit, fix $0<r<1$ and
choose
\[
\beta=\frac{1-r}{2},\qquad \alpha=\frac{2r}{1+r}.
\]
Then $\alpha(1-\beta)=r$, and the matrix reduces to
\begin{equation}\label{eq:P}
P_r=\begin{pmatrix}a&b\\b&a\end{pmatrix},\qquad
a=\frac{1+r}{2},\quad b=\frac{1-r}{2}.
\end{equation}
Here ``preserves the uniform law'' means that a site with input
probabilities $(1/2,1/2)$ has the same output probabilities:
\[
(\tfrac12,\tfrac12)P_r=(\tfrac12,\tfrac12).
\]
For functions on the two states, the matrix acts by conditional
expectation, $(P_rf)(s)=\sum_{u=0}^1P_r(s,u)f(u)$.
The function $\varepsilon(s)=(-1)^s$ has values $1,-1$, hence mean zero
under the uniform law; this is the meaning of ``centered.'' Directly,
\[
P_r\begin{pmatrix}1\\-1\end{pmatrix}
=\begin{pmatrix}a-b\\b-a\end{pmatrix}
=r\begin{pmatrix}1\\-1\end{pmatrix}.
\]
Thus $\mathbb E[\varepsilon(S_{\mathrm{out}})\mid S_{\mathrm{in}}=s]
=r\varepsilon(s)$: the constant mode has eigenvalue $1$, and the
mean-zero mode has eigenvalue $r$.

Apply $P_r$ independently to $N$ sites. Given $k$ input ones, let
$T_{r,N}(k,\ell)$ be the probability of $\ell$ output ones. Forgetting
site labels is the degree-$N$ induced-matrix construction. It gives
\begin{equation}\label{eq:transitionGF}
\sum_{\ell=0}^NT_{r,N}(k,\ell)z^\ell
=(a+bz)^{N-k}(b+az)^k.
\end{equation}
The stationary occupation law is the binomial law $\pi_N$. We define
$m_{r,N}$ to be the transition density \emph{relative to this law},
rather than the transition probability itself:
\begin{equation}\label{eq:density}
\pi_N(k)=2^{-N}\binom Nk,\qquad
m_{r,N}(k,\ell)=\frac{T_{r,N}(k,\ell)}{\pi_N(\ell)}.
\end{equation}
Thus $T_{r,N}(k,\ell)=m_{r,N}(k,\ell)\pi_N(\ell)$ and
$\sum_\ell m_{r,N}(k,\ell)\pi_N(\ell)=1$. With stationary input, it is
also the joint occupation law divided by the product of its margins:
\[
m_{r,N}(k,\ell)=
\frac{\Pr(K_{\mathrm{in}}=k,K_{\mathrm{out}}=\ell)}
{\pi_N(k)\pi_N(\ell)}.
\]
With stationary input, the probabilities of site pairs $00,01,10,11$
are $a/2,b/2,b/2,a/2$. Thus the joint occupation generating function is
\begin{equation}\label{eq:joint}
\sum_{k,\ell}\pi_N(k)T_{r,N}(k,\ell)u^kv^\ell
=\left[\frac{(1+u)(1+v)+r(1-u)(1-v)}4\right]^N.
\end{equation}
This is all the probabilistic structure required below.

\subsection{Kravchuk polynomials and the exact bilinear sum}
Define the ordinary-coefficient polynomials by
\begin{equation}\label{eq:Kdef}
(1+z)^{N-k}(1-z)^k=\sum_{n=0}^NK_n(k;N)z^n.
\end{equation}
For independent uniform input bits, the product of two such generating
functions has expectation $(1+zs)^N$. Comparing coefficients proves
\begin{equation}\label{eq:orthog}
\sum_k\pi_N(k)K_n(k;N)K_j(k;N)=\binom Nn\delta_{nj}.
\end{equation}
There is also the elementary transform
\[
\sum_k\pi_N(k)u^kK_n(k;N)
=\binom Nn\left(\frac{1+u}{2}\right)^{N-n}
\left(\frac{1-u}{2}\right)^n.
\]
Apply it in both variables to
$\sum_nr^nK_n(k;N)K_n(\ell;N)/\binom Nn$.
The binomial theorem gives exactly \eqref{eq:joint}.
Recall explicitly that the kernel on the left below is
\[
m_{r,N}(k,\ell):=\frac{T_{r,N}(k,\ell)}{2^{-N}\binom N\ell},
\]
where $T_{r,N}(k,\ell)$ is the output-count probability given $k$ input
successes. Equality of coefficients therefore proves the bilinear
identity
\begin{equation}\label{eq:finite-sum}
\boxed{\displaystyle
m_{r,N}(k,\ell)=\sum_{n=0}^N
\frac{r^n}{\binom Nn}K_n(k;N)K_n(\ell;N).}
\end{equation}
The model identifies the normalization and positivity; it does not yet
supply the scaling limit.

\section{The hypergeometric bilinear formula}
\subsection{The general identity and its finite specialization}
For $t\notin\{0,-1,-2,\ldots\}$ the bilinear identity is
\begin{equation}\label{eq:general-bilinear}
\begin{aligned}
&\sum_{n\ge0}\frac{(t)_n}{n!}w^n
{}_2F_1(-n,X;t;v){}_2F_1(-n,Y;t;v)\\
&\quad=(1-w)^{X+Y-t}A^{-X-Y}
{}_2F_1\left(X,Y;t;\frac{wv^2}{A^2}\right),\qquad
A=1-(1-v)w.
\end{aligned}
\end{equation}
It is first an identity near $w=0$, with the branches taking value $1$
there; Appendix~\ref{app:bilinear} gives a short proof.

To put $t=-N$, first set $X=-k$, $Y=-\ell$. Both terminating
polynomials then involve denominator factors only up to
$\min(k,\ell)\le N$. Passage to $t=-N$ is legitimate: in the series on
the left all $n>N$ terms vanish, while the right-hand series terminates
before a denominator zero. More explicitly, each factor on the left
is a polynomial in $n$ of bounded degree with coefficients analytic
near $t=-N$, so convergence is locally uniform for $|w|<1$.
Since $(-N)_n/n!=(-1)^n\binom Nn$, we obtain
\begin{equation}\label{eq:finite-v}
\begin{aligned}
&\sum_{n=0}^N(-1)^n\binom Nn w^n
{}_2F_1(-n,-k;-N;v){}_2F_1(-n,-\ell;-N;v)\\
&\quad=(1-w)^{N-k-\ell}A^{k+\ell}
{}_2F_1\left(-k,-\ell;-N;\frac{wv^2}{A^2}\right).
\end{aligned}
\end{equation}
If $k+\ell>N$, apparent denominators cancel. The right side is the
polynomial specified by the finite left side, including at removable
singularities.

\subsection{The closed finite kernel: the starting point for the limit}
In our convention,
\[
K_n(k;N)=\binom Nn{}_2F_1(-n,-k;-N;2).
\]
Take $v=2$, $w=-r$ in \eqref{eq:finite-v}. Equations
\eqref{eq:finite-sum} and \eqref{eq:finite-v} give
\begin{equation}\label{eq:closed-kernel}
\boxed{\begin{aligned}
m_{r,N}(k,\ell)
={}&(1+r)^{N-k-\ell}(1-r)^{k+\ell}\\
&\times{}_2F_1\left(-k,-\ell;-N;
-\frac{4r}{(1-r)^2}\right).
\end{aligned}}
\end{equation}
This closed bilinear expression is our asymptotic tool. Its parameters
$-k,-\ell,-N$ all grow, so fixed-parameter hypergeometric limits do not
apply. We evaluate it through its exact coefficient representation:
\begin{equation}\label{eq:coefficient}
m_{r,N}(k,\ell)
=\frac{[z^\ell](a+bz)^{N-k}(b+az)^k}
{2^{-N}\binom N\ell}.
\end{equation}
There is a direct algebraic conversion, independent of the model:
for nonzero $c$, the binomial theorem gives
\[
[z^\ell](1+cz)^{N-k}(1+dz)^k
=\binom N\ell c^\ell
{}_2F_1(-k,-\ell;-N;1-d/c).
\]
Set $c=b/a$, $d=a/b$ and restore the factor $a^{N-k}b^k$.
Then $1-d/c=-4r/(1-r)^2$, and division by
$2^{-N}\binom N\ell$ gives exactly \eqref{eq:closed-kernel}.
Thus \eqref{eq:coefficient} is an algebraic representation of the
closed bilinear formula. Its coefficient asymptotics below evaluate
that formula without assuming a joint central limit theorem.

\section{Asymptotics of the closed finite formula}
\begin{theorem}[Kravchuk bilinear kernel to Mehler]\label{thm:limit}
Fix $0<r<1$. For any sequences $k_N,\ell_N$ such that
$x_{k_N,N}\to x$ and $x_{\ell_N,N}\to y$,
\begin{equation}\label{eq:mehler-ratio}
\boxed{\displaystyle
m_{r,N}(k_N,\ell_N)\longrightarrow
m_r(x,y)=\frac1{\sqrt{1-r^2}}
\exp\left(\frac{2rxy-r^2(x^2+y^2)}{1-r^2}\right).}
\end{equation}
Convergence is uniform on compact central windows, and locally uniform
in $r\in(0,1)$.
\end{theorem}
\begin{proof}
We evaluate the numerator and denominator of \eqref{eq:coefficient}.
All error estimates below are uniform when the scaled coordinates
belong to a fixed compact set.

\smallskip\noindent\emph{Step 1: the stationary denominator.}
Stirling's formula in the central window gives
\begin{equation}\label{eq:binomial-asym}
2^{-N}\binom N\ell
=\sqrt{\frac2{\pi N}}\exp(-y_N^2)(1+o(1)),
\qquad y_N=x_{\ell,N}.
\end{equation}
The prefactor is $\delta_Ng(y_N)$, as required by the lattice spacing.

\smallskip\noindent\emph{Step 2: a coefficient integral.}
Put $x_N=x_{k,N}$ and write the numerator as
\begin{equation}\label{eq:integral}
C_N=\frac1{2\pi}\int_{-\pi}^{\pi}
e^{-i\ell\theta}(a+be^{i\theta})^{N-k}
(b+ae^{i\theta})^k\,d\theta.
\end{equation}
Near $\theta=0$, logarithms continuous from $0$ satisfy
\[
\begin{aligned}
\log(a+be^{i\theta})&=ib\theta-\tfrac12ab\theta^2+O(|\theta|^3),\\
\log(b+ae^{i\theta})&=ia\theta-\tfrac12ab\theta^2+O(|\theta|^3).
\end{aligned}
\]
Their combined linear coefficient before coefficient extraction is
\[
\mu_N=(N-k)b+ka=\frac N2+r\left(k-\frac N2\right),
\qquad
\mu_N-\ell=\sqrt{\frac N2}(y_N-rx_N).
\]
Consequently, for $\theta=u/\sqrt N$, the exponent in the integrand is
\begin{equation}\label{eq:local-exponent}
\frac{i(y_N-rx_N)u}{\sqrt2}
-\frac{1-r^2}{8}u^2+O\left(\frac{|u|^3}{\sqrt N}\right).
\end{equation}
This is a direct saddle-point expansion of the coefficient form of
\eqref{eq:closed-kernel}.

\smallskip\noindent\emph{Step 3: justify localization.}
Both factors have the same modulus squared:
\[
|a+be^{i\theta}|^2=|b+ae^{i\theta}|^2
=1-(1-r^2)\sin^2(\theta/2).
\]
Thus the modulus of their product in \eqref{eq:integral} is at most
\[
\exp\left[-\frac N2(1-r^2)\sin^2(\theta/2)\right].
\]
Since $\sin(|\theta|/2)\ge |\theta|/\pi$ on $[-\pi,\pi]$, the scaled
integrand is bounded by the integrable function
$\exp[-(1-r^2)u^2/(2\pi^2)]$, independently of $N,k,\ell$.
A fixed-$u$ expansion followed by dominated convergence therefore
justifies the Gaussian integral, without using a probabilistic limit
theorem. Uniformity follows by the same bound and uniform expansion
on compact sets.

\smallskip\noindent\emph{Step 4: evaluate and divide.}
Equations \eqref{eq:integral}--\eqref{eq:local-exponent} give
\begin{equation}\label{eq:numerator-asym}
C_N=\sqrt{\frac{2}{\pi N(1-r^2)}}
\exp\left[-\frac{(y_N-rx_N)^2}{1-r^2}\right](1+o(1)).
\end{equation}
Dividing by \eqref{eq:binomial-asym},
\[
m_{r,N}(k,\ell)
=\frac1{\sqrt{1-r^2}}
\exp\left[y_N^2-\frac{(y_N-rx_N)^2}{1-r^2}\right](1+o(1)).
\]
The exponent equals
$[2rx_Ny_N-r^2(x_N^2+y_N^2)]/(1-r^2)$, proving the theorem.
\end{proof}

Combining the theorem with the closed bilinear formula makes the
hypergeometric asymptotic explicit:
\begin{equation}\label{eq:hypergeo-limit}
\begin{aligned}
&(1+r)^{N-k_N-\ell_N}(1-r)^{k_N+\ell_N}
{}_2F_1\left(-k_N,-\ell_N;-N;-\frac{4r}{(1-r)^2}\right)\\
&\hspace{2em}\longrightarrow
\frac1{\sqrt{1-r^2}}
\exp\left(\frac{2rxy-r^2(x^2+y^2)}{1-r^2}\right).
\end{aligned}
\end{equation}
The coefficient representation has handled the large cancellations in
the terminating hypergeometric polynomial. Expanding its summands
individually would not justify this limit.

\section{Identification as the Mehler kernel}
\subsection{Density relative to Gaussian measure}
Multiplying \eqref{eq:mehler-ratio} by $g(y)$ gives
\begin{equation}\label{eq:gaussian-transition}
m_r(x,y)g(y)=\frac1{\sqrt{\pi(1-r^2)}}
\exp\left[-\frac{(y-rx)^2}{1-r^2}\right].
\end{equation}
It integrates to $1$. Thus the limiting bilinear expression is the
Mehler transition density relative to $d\mu(y)=g(y)dy$, not relative
to Lebesgue measure. This distinction accounts for the quadratic terms
in \eqref{eq:mehler-ratio}.

\subsection{Hermite normalization and the bilinear series}
Let $H_n$ denote the physicists' Hermite polynomials,
\[
e^{2xs-s^2}=\sum_{n\ge0}H_n(x)\frac{s^n}{n!},\qquad
\int H_nH_j\,d\mu=2^nn!\delta_{nj}.
\]
For fixed $n$, rescale \eqref{eq:Kdef} by $z=\sqrt{2/N}\,s$.
Its logarithm tends locally uniformly to $2xs-s^2$, so
\begin{equation}\label{eq:fixed-degree}
\frac{K_n(k_N;N)}{\sqrt{\binom Nn}}
\longrightarrow\frac{H_n(x)}{\sqrt{2^nn!}}.
\end{equation}
This identifies the individual modes but was not used to interchange
the finite sum with the limit.

For completeness, Gaussian integration in \eqref{eq:gaussian-transition}
gives
\[
\int m_r(x,y)e^{2ys-s^2}\,d\mu(y)=e^{2rxs-r^2s^2}.
\]
Comparing coefficients shows that $H_n$ has multiplier $r^n$.
Orthogonality and completeness therefore give the Mehler expansion
\begin{equation}\label{eq:hermite-mehler}
\boxed{\displaystyle
\sum_{n=0}^{\infty}\frac{r^n}{2^nn!}H_n(x)H_n(y)
=\frac1{\sqrt{1-r^2}}
\exp\left(\frac{2rxy-r^2(x^2+y^2)}{1-r^2}\right).}
\end{equation}
The expansion is absolutely convergent for $|r|<1$ and real $x,y$;
Cauchy estimates for the Hermite generating function give subexponential
growth of the normalized modes on compact $x$-sets. In particular the
series identifies the continuous kernel, not merely its $L^2$ class.
Our limit proof concerns $0<r<1$ and has not assumed this expansion.

\subsection{The Hermite-function version}
If a kernel with respect to $dy$ rather than $d\mu(y)$ is desired,
symmetric Gaussian conjugation gives
\begin{equation}\label{eq:wavefunction-kernel}
\begin{aligned}
\sqrt{g(x)}m_r(x,y)\sqrt{g(y)}
={}&\frac1{\sqrt{\pi(1-r^2)}}\\
&\times\exp\left[-\frac{(1+r^2)(x^2+y^2)-4rxy}{2(1-r^2)}\right].
\end{aligned}
\end{equation}
Equivalently this is $\sum_{n\ge0}r^n\phi_n(x)\phi_n(y)$ for
$\phi_n(x)=\pi^{-1/4}e^{-x^2/2}H_n(x)/\sqrt{2^nn!}$.
This is only a normalization change of the Mehler kernel already
obtained; no additional operator theory is needed.

\appendix
\section{Proof of the general bilinear identity}\label{app:bilinear}
Initially take $t>X,Y>0$, $0<v<1$, and $w>0$ small. Euler's beta
integral represents ${}_2F_1(-n,X;t;v)$ as the average of $(1-vs)^n$
against the beta$(X,t-X)$ density. Use it for both factors in
\eqref{eq:general-bilinear} and sum
$\sum_n(t)_nz^n/n!=(1-z)^{-t}$. The resulting double integral has
integrand $[1-w(1-vs)(1-vu)]^{-t}$.

Integrating over $s$ first, and setting $c=1-w$, $d=wv$, $A=c+d$,
gives
\[
\frac1{B(Y,t-Y)}\int_0^1
u^{Y-1}(1-u)^{t-Y-1}
(c+du)^{X-t}[A+d(1-v)u]^{-X}\,du.
\]
Make the substitution $u=ch/(A-dh)$. The integral becomes
\[
\frac{c^{X+Y-t}A^{-X-Y}}{B(Y,t-Y)}\int_0^1
h^{Y-1}(1-h)^{t-Y-1}
\left(1-\frac{wv^2}{A^2}h\right)^{-X}dh.
\]
Euler's integral now gives \eqref{eq:general-bilinear}. Analytic
continuation proves it for general parameters wherever the indicated
functions are defined, initially as analytic germs at $w=0$.


\begin{thebibliography}{9}
\bibitem{FeinsilverSFAII2023}
P.~Feinsilver, \emph{Symmetric Functions Algebras (SFA) II: Induced
Matrices}, Journal of Stochastic Analysis \textbf{4} (2023), no.~2,
Article~1.
\href{https://doi.org/10.31390/josa.4.2.01}{doi:10.31390/josa.4.2.01}.

\bibitem{HoareRahman1984}
M.~R. Hoare and M.~Rahman, \emph{Cumulative Bernoulli trials and
Krawtchouk processes}, Stochastic Process. Appl. \textbf{16} (1984),
no.~2, 113--139.
\href{https://doi.org/10.1016/0304-4149(84)90014-0}{doi:10.1016/0304-4149(84)90014-0}.
The publisher dates the issue February 1984; its copyright notice is 1983.

\end{thebibliography}
\end{document}
